A ground robot that loses satellite positioning in a desert must keep its heading while its cameras degrade in dust, glare and low light. Desert ants use a sky polarisation compass for this. We test, in a third and stronger study, when such a compass helps stereo visual odometry and how a filter should weight it. A vector Monte Carlo radiative transfer sky with climatological aerosol, checked against measured desert values and one day of open full sky polarisation images, brackets the noon blind window of a zenith compass in Abu Dhabi on 21 June at \sWinAnch{} to \sWinRT{} minutes and places such windows at every latitude within about \sLatMax{} degrees of the equator. A Cram\'er Rao analysis shows that no sky geometry removes the window, while a ring of off zenith points shortens it to minutes. With a compass bias, exact tilt geometry, dust tied to the image visibility and a filter that no longer knows the true compass noise, we compared nineteen systems under a protocol frozen before a confirmation run on \sNconf{} rendered trajectories from three environments never used before. A weighted compass lowered the heading RMSE from \sHoneB{} to \sHoneA{} degrees; an ungated compass was worse than none even with clear images; DoP floored covariance matching was non inferior to covariance matching; a ring geometry beat the zenith point; and simulated wheel odometry with a sky compass drifted \sHfiveA{} percent of path length at the most severe level against \sHfiveB{} for vision alone. A calibrated noise model did not repair the overconfidence, which an oracle analysis traces to biased visual translation, and the ranking among weighting schemes did not transfer between splits. Scenes are rendered and the compass is simulated.
